\chapter{Managing a Book through Drawdown}\label{ch:fm:managing-a-book-through-drawdown}

Two portfolio managers are both down eight per cent from their peak in March. One runs a strategy that has had a bad few months; the other's
edge died a year ago and nobody has noticed. On the day each first reaches the eight per cent, the best statistical reading of their P\&L gives
the first a 24\% chance of having broken and the second 39\%: the numbers barely differ. To tell a strategy whose Sharpe ratio has fallen from 1
to 0 from one that is merely unlucky, at a false-alarm rate of one in ten years, takes on average more than fifteen years of data. The head of desk
must decide in March.

\section{What a drawdown says about a Sharpe ratio}

A drawdown (One Quant Book 2, chapter 29) is the fall of cumulative P\&L from its running peak; its maximum over a period and its duration are
Book 7's measures (chapter 22). What matters for management is how large a drawdown a \emph{healthy} strategy produces. For daily P\&L with a
Sharpe ratio of 1 and a volatility of 10\% a year on capital, the maximum drawdown over ten years has a median of 16.8\% of capital and a
90th percentile of 25.4\% (2\,000 simulated paths): a healthy strategy of this kind will spend some years more than one year's volatility below
its peak. Magdon-Ismail, Atiya, Pratap and Abu-Mostafa give the distribution of the maximum drawdown of a Brownian motion with drift; the
simulation reproduces its lesson, that drawdowns of one to two years' volatility are the normal life of a Sharpe-1 strategy.

\begin{remark}[Units]\label{rem:fm:managing-a-book-through-drawdown:units}
A drawdown means nothing until it is divided by the strategy's volatility. Eight per cent is 0.8 of a year's volatility at 10\% and 0.27 at 30\%.
Every rule in this chapter is stated for a strategy with 10\% volatility; for another, scale the thresholds by the ratio of volatilities (One
Quant Book 8, chapter 28, reached the same conclusion from the other side).
\end{remark}

\section{Stop rules and de-risking ladders}

\begin{definition}[Stop-loss rule, de-risking ladder, time stop]\label{def:fm:managing-a-book-through-drawdown:rules}
A \emph{stop-loss rule}\index{stop-loss rule} stops a strategy when its drawdown from its running peak exceeds a threshold. A \emph{de-risking
ladder}\index{de-risking ladder} cuts a strategy's size in steps as its drawdown deepens (halving at a first threshold, stopping at a second)
instead of stopping it at once. A \emph{time stop}\index{time stop} stops a strategy that has made no new high for a given length of time,
whatever the depth of its drawdown.
\end{definition}

A stop rule makes two errors, and the drawdown limit of One Quant Book 8, chapter 28 already showed their shape: it stops live strategies in bad
luck, and it lets dead ones run. The ladder is a two-step version of Grossman and Zhou's rule, in which the position falls smoothly as the
drawdown approaches a floor; the \omterm{def:fm:managing-a-book-through-drawdown:rules}{time stop} watches duration instead of depth, because a dead strategy with a Sharpe ratio of zero drifts
sideways rather than down.

\omcode{code/firm/ddrules/firm_ddrules.py}{40}{55}{The posterior probability that a strategy has broken: a forward filter on a two-state chain,
alive and broken, with a daily hazard of breaking.}\label{lst:fm:managing-a-book-through-drawdown:posterior}

\section{Bad luck or broken: sequential evidence}

A drawdown rule looks at one statistic of the P\&L path. The whole path carries more: every day's return is evidence about whether the strategy
is still the one that was backtested.

\begin{proposition}[The posterior that a strategy has broken]\label{prop:fm:managing-a-book-through-drawdown:posterior}
Let a strategy be alive or broken; alive, its daily return is $\mathcal N(\mu_1,\sigma_d^2)$; broken, $\mathcal N(\mu_0,\sigma_d^2)$; it breaks on
any day with probability $h$ if not yet broken, and never recovers. With $p_{t-1}$ the probability that it is broken given the returns up to day
$t-1$, the probability given day $t$'s return $x_t$ is
\[
p_t=\frac{\tilde p_t\,\varphi\!\left(\frac{x_t-\mu_0}{\sigma_d}\right)}{\tilde p_t\,\varphi\!\left(\frac{x_t-\mu_0}{\sigma_d}\right)+(1-\tilde p_t)\,\varphi\!\left(\frac{x_t-\mu_1}{\sigma_d}\right)},
\qquad \tilde p_t=p_{t-1}+(1-p_{t-1})\,h .
\]
\end{proposition}

\begin{proof}
The prediction step: the strategy is broken on day $t$ if it was broken before or breaks today, $\tilde p_t=p_{t-1}+(1-p_{t-1})h$. The update is
Bayes' rule with the two Gaussian likelihoods; the common factor $1/\sigma_d$ cancels.
\end{proof}

\begin{proposition}[How long it takes to know]\label{prop:fm:managing-a-book-through-drawdown:lorden}
For the detection of a change from the alive to the broken distribution with an average of $A$ days between false alarms, no detector has, as
$A\to\infty$, a worst-case expected delay shorter than $\ln A/K$, and the CUSUM detector attains it, where $K$ is the Kullback--Leibler divergence
per day. For daily returns with annual Sharpe ratios $S$ before and $S_0$ after the break, $K=(S-S_0)^2/(2\times252)$, so the delay is
\[
\frac{2\,\ln A}{(S-S_0)^2}\ \text{years}.
\]
\end{proposition}

\admitted

The bound is Lorden's (1971) result, cited in \texttt{omsources}; the CUSUM test itself is One Quant Book 7's (chapter 13). Its arithmetic is the
chapter's central fact. A Sharpe ratio that falls from 1 to 0, detected at one false alarm per ten years ($A=2\,520$ days), takes
$2\ln2\,520=15.7$ years to detect on average; at one per five years, 14.3. A fall to $-1$ takes 3.9 years and a fall to $-2$, 1.7
(\cref{fig:fm:managing-a-book-through-drawdown:lorden}). An edge that disappears quietly cannot be distinguished from bad luck on its P\&L in any
useful time; only an edge that turns into a loss can.

\begin{omfigure}
\begin{tikzpicture}
\begin{axis}[width=10.5cm,height=5.2cm,xlabel={\small Sharpe ratio after the break $S_0$ (before: 1)},ylabel={\small years to detect},
  xmin=-2,xmax=0.5,ymin=0,ymax=65,tick label style={font=\footnotesize},grid=major,grid style={gray!20},table/col sep=comma]
\addplot[omDef,thick] table[x=sr_dead,y=years]{figdata/desk/08-managing-a-book-through-drawdown/lorden.csv};
\draw[omThm,dashed] (axis cs:0,0) -- (axis cs:0,65);
\node[font=\scriptsize,anchor=west,fill=white,inner sep=1pt] at (axis cs:-1.9,50) {one false alarm per ten years};
\end{axis}
\end{tikzpicture}

\omcaption{The shortest average delay with which any detector can tell that a strategy's Sharpe ratio has fallen from 1 to $S_0$, at one false
alarm per ten years, from \cref{prop:fm:managing-a-book-through-drawdown:lorden}: 15.7 years for a fall to zero (dashed line), 3.9 for a fall to
$-1$. Data: \texttt{fm\_drawdown.lorden}.}\label{fig:fm:managing-a-book-through-drawdown:lorden}
\end{omfigure}

The posterior of \cref{prop:fm:managing-a-book-through-drawdown:posterior}, run on two example paths, shows what the head of desk sees
(\cref{fig:fm:managing-a-book-through-drawdown:paths}). The live strategy's posterior wanders, reaching 0.80 in its fourth year; the broken
strategy's climbs slowly after its break at the start of year four and passes 0.9 only in its fifth year.

\begin{omfigure}
\begin{tikzpicture}
\begin{axis}[width=10.5cm,height=4.4cm,ylabel={\small cumulative P\&L (\%)},xmin=0,xmax=10,ymin=-20,ymax=100,
  tick label style={font=\footnotesize},grid=major,grid style={gray!20},table/col sep=comma,xticklabels={},
  legend cell align=left,legend style={font=\scriptsize,at={(0.02,0.97)},anchor=north west}]
\addplot[omDef,thick] table[x=year,y=eq_alive]{figdata/desk/08-managing-a-book-through-drawdown/examples.csv};
\addplot[omThm,thick,dashed] table[x=year,y=eq_broken]{figdata/desk/08-managing-a-book-through-drawdown/examples.csv};
\draw[gray] (axis cs:3,-20) -- (axis cs:3,100);
\legend{alive,{breaks at year 3}}
\end{axis}
\begin{axis}[width=10.5cm,height=4.4cm,yshift=-3.6cm,xlabel={\small year},ylabel={\small P(broken)},xmin=0,xmax=10,ymin=0,ymax=1.05,
  tick label style={font=\footnotesize},grid=major,grid style={gray!20},table/col sep=comma]
\addplot[omDef,thick] table[x=year,y=post_alive]{figdata/desk/08-managing-a-book-through-drawdown/examples.csv};
\addplot[omThm,thick,dashed] table[x=year,y=post_broken]{figdata/desk/08-managing-a-book-through-drawdown/examples.csv};
\draw[gray] (axis cs:3,0) -- (axis cs:3,1.05);
\draw[omProp,dotted,thick] (axis cs:0,0.9) -- (axis cs:10,0.9);
\end{axis}
\end{tikzpicture}

\omcaption{Two strategies with a Sharpe ratio of 1 and 10\% volatility; the second's Sharpe ratio falls to 0 at the start of year 4 (vertical
line). Top: cumulative P\&L. Bottom: the posterior probability that each has broken (\cref{prop:fm:managing-a-book-through-drawdown:posterior},
hazard once in five years), with the 0.9 threshold dotted. Data: \texttt{fm\_drawdown.examples}.}\label{fig:fm:managing-a-book-through-drawdown:paths}
\end{omfigure}

\section{Comparing the rules}

The chapter's evaluation runs 2\,000 live strategies and 2\,000 that break at the start of year 4 to a Sharpe ratio of 0, for ten years, under each
rule, with every stop permanent. It reports false stops per hundred live strategy-years, the share of broken strategies stopped after their break
and the median delay, and the value kept per strategy over ten years: P\&L less a \omterm{def:fm:limits-and-capital-allocation:raroc}{capital charge} of 2\% of capital for each year run at full
size (a 15\% hurdle on capital of 13\% of the strategy's size), in per cent of capital.

\begin{center}
\small
\begin{tabular}{@{}lrrrrrr@{}}
\toprule
rule & false stops & broken & median & value, & value, & value,\\
 & per 100 years & caught & delay (years) & live & broken & half each\\
\midrule
none & -- & -- & -- & 79.0 & 11.1 & 45.1\\
stop-loss 10\% & 9.79 & 34\% & 0.72 & 21.7 & 15.0 & 18.3\\
ladder 7.5 / 12.5\% & 2.09 & 60\% & 2.79 & 39.5 & 12.8 & 26.2\\
\omterm{def:fm:managing-a-book-through-drawdown:rules}{time stop} 18 months & 6.19 & 80\% & 2.00 & 54.6 & 19.5 & 37.1\\
posterior above 0.9 & 2.84 & 95\% & 2.68 & 71.4 & 18.2 & 44.8\\
\bottomrule
\end{tabular}
\end{center}

\begin{omfigure}
\begin{tikzpicture}
\begin{axis}[width=11cm,height=5.2cm,ybar,bar width=8pt,ymin=0,ymax=85,ylabel={\small value kept (\% of capital)},
  xtick={0,1,2,3,4},xticklabels={none,{stop-loss 10\%},{ladder 7.5 / 12.5\%},{time stop 18 months},{posterior $>0.9$}},
  xticklabel style={font=\scriptsize,text width=1.8cm,align=center},tick label style={font=\footnotesize},table/col sep=comma,
  enlarge x limits=0.12,grid=major,grid style={gray!20},legend cell align=left,
  legend style={font=\scriptsize,at={(0.5,-0.32)},anchor=north,/tikz/every even column/.append style={column sep=8pt}},legend columns=3]
\addplot[fill=omDef!55,draw=omDef] table[x=k,y=value_alive]{figdata/desk/08-managing-a-book-through-drawdown/rules.csv};
\addplot[fill=omThm!50,draw=omThm] table[x=k,y=value_dead]{figdata/desk/08-managing-a-book-through-drawdown/rules.csv};
\addplot[fill=gray!40,draw=gray] table[x=k,y=value_mix]{figdata/desk/08-managing-a-book-through-drawdown/rules.csv};
\legend{live strategies,broken strategies,half of each}
\end{axis}
\end{tikzpicture}

\omcaption{Value kept per strategy over ten years (P\&L less a 2\%-a-year \omterm{def:fm:limits-and-capital-allocation:raroc}{capital charge} while running, in per cent of capital) under each rule,
for strategies that stay alive, strategies that break at year 3 to a Sharpe ratio of 0, and a population of half each. 2\,000 paths of each; stops
are permanent. Data: \texttt{fm\_drawdown.table}.}\label{fig:fm:managing-a-book-through-drawdown:rules}
\end{omfigure}

The 10\% stop-loss is the worst rule by far. It stops 98\% of healthy strategies at some point in ten years, and two thirds of the broken ones
\emph{before} they break; it keeps 18.3\% of capital per strategy against 45.1\% for no rule at all. The posterior rule stops 28\% of healthy
strategies over ten years, catches 95\% of the broken ones about two and a half years after the break, and keeps 44.8\%, as much as doing
nothing. Stopping a broken strategy whose Sharpe ratio is zero saves no P\&L in expectation; it saves the capital it ties up and the risk it adds,
which is why the \omterm{def:fm:limits-and-capital-allocation:raroc}{capital charge} is in the value. With a charge of zero no rule beats doing nothing; with a strategy that turns into a loser, every
rule catches it sooner and the posterior rule is clearly best (exercise 7).

\begin{method}[Reviewing a strategy in drawdown]\label{met:fm:managing-a-book-through-drawdown:review}
\begin{enumerate}
\item Express the drawdown and its duration in units of the strategy's volatility; compare them with the distribution for a healthy strategy.
\item Update the posterior that it has broken from its whole P\&L history, with a hazard set from the firm's experience of how long strategies
live (One Quant Book 11, chapter 28).
\item Look for evidence that is not P\&L: capture and fill rates, crowding measures (One Quant Book 7, chapter 28), a change in the market's
structure. These separate bad luck from broken far faster than P\&L does.
\item Cut size gradually as evidence accumulates; stop only on the posterior or on non-P\&L evidence; never on depth alone.
\end{enumerate}
\end{method}

\section{Restart, re-underwrite, retire}

\begin{definition}[Restart rule]\label{def:fm:managing-a-book-through-drawdown:restart}
A \emph{restart rule}\index{restart rule} states when and at what size a stopped strategy may trade again: after a fixed period, after its paper
P\&L recovers, or after a review of its thesis finds it intact.
\end{definition}

A stop is a decision under uncertainty; a \omterm{def:fm:managing-a-book-through-drawdown:restart}{restart rule} is what makes it reversible. Because a healthy strategy is stopped by any rule that
catches broken ones in useful time, the cost of false stops is paid in the months the strategy is flat. A \omterm{def:fm:managing-a-book-through-drawdown:restart}{restart rule} that paper-trades a stopped
strategy and brings it back at reduced size when its paper P\&L recovers recovers much of that cost; a restart that happens automatically at the
new year, as in the model of One Quant Book 8, chapter 28, recovers it only for the good strategies and restarts the dead ones too.
\emph{Re-underwriting} a strategy is reviewing it as if it were new: its thesis, its evidence, its capacity and its cost of capital. A
strategy that fails that review is retired even if its P\&L has not yet said so, the only way to beat the delays of
\cref{prop:fm:managing-a-book-through-drawdown:lorden}.

\begin{remark}[The people in a drawdown]\label{rem:fm:managing-a-book-through-drawdown:people}
A portfolio manager in drawdown faces rules that are public and asymmetric: a stop ends the job, a recovery restores the payout only above the
high-water mark (chapter 3). Both push towards taking more risk to recover quickly or less to avoid the stop. A ladder that cuts size before the
stop, and a \omterm{def:fm:managing-a-book-through-drawdown:restart}{restart rule} that does not end a career at the first stop, reduce the incentive to gamble for resurrection; so does a firm that is
known to review theses, not only P\&L.
\end{remark}

\section{Tutorial: two thousand strategies and four rules}\label{tut:fm:managing-a-book-through-drawdown:tut}

\begin{tutorial}
\textbf{Goal.} Measure what each drawdown rule costs a healthy strategy and saves on a broken one. \textbf{End state:} the table of rules and
\cref{fig:fm:managing-a-book-through-drawdown:rules}.
\begin{tutsteps}
\item \textbf{Paths.} \texttt{fm\_drawdown.population()} draws 2\,000 live and 2\,000 breaking strategies with \texttt{firm.ddrules.paths}.
\item \textbf{Posterior.} \texttt{firm.ddrules.posterior} runs the filter of \cref{lst:fm:managing-a-book-through-drawdown:posterior} on every
path.
\item \textbf{Rules.} \texttt{fm\_drawdown.table()} applies the four rules with \texttt{apply} and scores them with \texttt{evaluate}.
\item \textbf{The bound.} \texttt{fm\_drawdown.lorden(10, S0)} evaluates \cref{prop:fm:managing-a-book-through-drawdown:lorden}.
\end{tutsteps}
\end{tutorial}

\textbf{What to change next.} Break the strategies to a Sharpe ratio of $-1$ (exercise 7); set the hazard ten times too high and watch the
posterior rule stop healthy strategies.

\section{Build: drawdown rules}\label{bld:fm:managing-a-book-through-drawdown:ddrules}

\begin{build}
\bfield{Purpose} Rules for managing strategies in drawdown, and an honest measure of what each costs and saves.
\bfield{Interface} \texttt{firm.ddrules}: \texttt{paths(n, days, sr, vol, rng, break\_day, sr\_dead)}; \texttt{posterior(returns, sr, sr\_dead,
vol, hazard)}; \texttt{Rule(kind, a, b)} with kinds \texttt{stop}, \texttt{ladder}, \texttt{time}, \texttt{posterior}; \texttt{apply};
\texttt{evaluate(rule, alive, dead, break\_day, post\_alive, post\_dead, charge)}; \texttt{cusum\_delay(sr, sr\_dead, arl\_days)}.
\bfield{Rules} A day's size is decided on the P\&L up to the day before; stops are permanent in the evaluator; value deducts the \omterm{def:fm:limits-and-capital-allocation:raroc}{capital charge}
only while the strategy runs.
\bfield{Acceptance tests} \texttt{code/firm/ddrules/tests/}: each rule on a hand-made path (the stop day, the halving, the \omterm{def:fm:managing-a-book-through-drawdown:rules}{time stop}); the
posterior separates live from broken populations; the evaluator's shape; the delay formula.
\bfield{Stretch} \omterm{def:fm:managing-a-book-through-drawdown:restart}{Restart rules} in the evaluator; a posterior over the post-break Sharpe ratio instead of a fixed $S_0$; evidence from fill rates
and capture as a second observation stream.
\end{build}

\begin{omsources}
\item M.~Magdon-Ismail, A.~F.~Atiya, A.~Pratap and Y.~S.~Abu-Mostafa, ``On the maximum drawdown of a Brownian motion'', \emph{Journal of
Applied Probability} 41(1), 2004.
\item G.~Lorden, ``Procedures for reacting to a change in distribution'', \emph{Annals of Mathematical Statistics} 42(6), 1971.
\item R.~P.~Adams and D.~J.~C.~MacKay, ``Bayesian online changepoint detection'', arXiv:0710.3742, 2007.
\end{omsources}

\section{Exercises}

\begin{exercise}[$\star$]\label{exo:fm:managing-a-book-through-drawdown:1}
Express a drawdown of 8\% in years of volatility for strategies with volatilities of 10\% and 30\%.
\end{exercise}

\begin{exercise}[$\star$]\label{exo:fm:managing-a-book-through-drawdown:2}
Use \cref{prop:fm:managing-a-book-through-drawdown:lorden} to compute the delay to detect a fall of the Sharpe ratio from 1 to 0 at one false alarm
per ten years, and from 1 to $-1$.
\end{exercise}

\begin{exercise}[$\star$]\label{exo:fm:managing-a-book-through-drawdown:3}
From the table of rules, what share of healthy strategies does the 10\% stop-loss stop over ten years, and what share of broken ones does it stop
before they break?
\end{exercise}

\begin{exercise}[$\star\star$]\label{exo:fm:managing-a-book-through-drawdown:4}
A strategy's posterior of having broken is 0.30 and the daily hazard is $1/1\,260$. Its next day's return is exactly the broken mean. Is the
posterior higher or lower after that day, and why does one day move it so little?
\end{exercise}

\begin{exercise}[$\star\star$]\label{exo:fm:managing-a-book-through-drawdown:5}
Why does stopping a broken strategy whose Sharpe ratio is zero add no expected P\&L, and what does it save?
\end{exercise}

\begin{exercise}[$\star\star$]\label{exo:fm:managing-a-book-through-drawdown:6}
A strategy has a Sharpe ratio of 1 and a volatility of 10\%. Its maximum drawdown over ten years has a median of 16.8\%. Is a 12\% drawdown in its
sixth year evidence that it has broken?
\end{exercise}

\begin{exercise}[$\star\star\star$]\label{exo:fm:managing-a-book-through-drawdown:7}
\emph{Coding.} Rerun the evaluation with broken strategies falling to a Sharpe ratio of $-1$ (set \texttt{SR\_DEAD} and the posterior's
alternative to $-1$). What happens to the delays and to each rule's value?
\end{exercise}

\begin{exercise}[$\star\star\star$]\label{exo:fm:managing-a-book-through-drawdown:8}
\emph{Find the flaw.} ``We backtested our 10\% stop-loss on our strategies' histories and it would have avoided every large loss. Adopt it.''
\end{exercise}

\section{Problem: Down Eight in March}

\begin{problem}[{Weekend problem --- down eight in March}]\label{pb:fm:managing-a-book-through-drawdown:1}
Two portfolio managers are each down 8\% from their peak. The head of desk must decide what to do with each, and what rule to adopt for next time.

\medskip\noindent\textbf{Part I --- The drawdown.}
\begin{enumerate}
\item Define a \omterm{def:fm:managing-a-book-through-drawdown:rules}{stop-loss rule}, a \omterm{def:fm:managing-a-book-through-drawdown:rules}{de-risking ladder} and a \omterm{def:fm:managing-a-book-through-drawdown:rules}{time stop}.
\item Give the median and 90th percentile of a healthy Sharpe-1 strategy's maximum drawdown over ten years at 10\% volatility.
\item Express 8\% in volatility units for both managers, at 10\% and 20\% volatility.
\item Why do thresholds belong in volatility units?
\end{enumerate}

\medskip\noindent\textbf{Part II --- The evidence.}
\begin{enumerate}[resume]
\item State and prove \cref{prop:fm:managing-a-book-through-drawdown:posterior}.
\item Give the mean posterior on the day a healthy and a broken strategy first reach an 8\% drawdown.
\item State \cref{prop:fm:managing-a-book-through-drawdown:lorden} and compute the delay for a fall from 1 to 0 and to $-1$ at one false alarm per
ten years.
\item Which kinds of evidence separate bad luck from broken faster than P\&L?
\end{enumerate}

\medskip\noindent\textbf{Part III --- The rules.}
\begin{enumerate}[resume]
\item Give each rule's false stops per hundred strategy-years and the share of broken strategies caught after their break.
\item Give each rule's median detection delay.
\item Give each rule's value kept for live, broken and mixed populations, against no rule.
\item Why is the stop-loss the worst rule?
\item Why does the posterior rule keep about as much as no rule, and when would it keep more?
\end{enumerate}

\medskip\noindent\textbf{Part IV --- The decision.}
\begin{enumerate}[resume]
\item Define a \omterm{def:fm:managing-a-book-through-drawdown:restart}{restart rule} and describe a good one.
\item What is re-underwriting, and why can it beat the delays of the P\&L?
\item What do the rules do to a portfolio manager's incentives in drawdown?
\item What would you do with each of the two managers on Monday?
\item Which rule would you adopt, with which thresholds, and why?
\item State the \emph{named result}: the posterior on the day of the 8\% drawdown for a healthy and a broken strategy, and the detection delay of
a fall from 1 to 0 at one false alarm per ten years.
\item In two sentences, explain to the managers how they will be judged.
\end{enumerate}
\end{problem}

\section{Interview questions}

\begin{interviewq}[$\star$ \iqroles{trader}]\label{iq:fm:managing-a-book-through-drawdown:1}
Your strategy has a Sharpe ratio of 1 and 10\% volatility and is 15\% below its peak. Should you be worried?
\end{interviewq}

\begin{interviewq}[$\star$ \iqroles{risk}]\label{iq:fm:managing-a-book-through-drawdown:2}
What is the difference between a stop-loss and a \omterm{def:fm:managing-a-book-through-drawdown:rules}{de-risking ladder}, and why might a firm prefer the ladder?
\end{interviewq}

\begin{interviewq}[$\star\star$ \iqroles{researcher}]\label{iq:fm:managing-a-book-through-drawdown:3}
How long would it take to detect, from daily P\&L, that a strategy's Sharpe ratio has fallen from 1 to 0? Derive the order of magnitude.
\end{interviewq}

\begin{interviewq}[$\star\star$ \iqroles{researcher, risk}]\label{iq:fm:managing-a-book-through-drawdown:4}
Write the update of the probability that a strategy has broken, given one more day of returns.
\end{interviewq}

\begin{interviewq}[$\star\star$ \iqroles{trader, risk}]\label{iq:fm:managing-a-book-through-drawdown:5}
A backtest shows that a stop-loss would have avoided every large loss. What would you check?
\end{interviewq}

\begin{interviewq}[$\star\star\star$ \iqroles{researcher}]\label{iq:fm:managing-a-book-through-drawdown:6}
Your P\&L cannot tell you in time whether an edge has gone. What else would you monitor, and why is it faster?
\end{interviewq}
